\documentclass[10pt,a4paper]{amsart}
\usepackage[margin=19mm]{geometry}
\usepackage{amsmath,amssymb,mathtools}
\usepackage[scr=rsfs,cal=euler]{mathalfa}
\usepackage{xfrac}
\usepackage[unicode,hidelinks,bookmarks=false]{hyperref}
\newcommand{\ie}{i.e.\ }
\newcommand{\cf}{cf.\ }
\newcommand{\resp}{respectively\ }
\newcommand{\Subsection}{\subsection}
\providecommand{\nobreakdash}{\nobreak\hbox{-}}
\setlength{\emergencystretch}{2em}
\multlinegap0pt
\usepackage{bm}
\usepackage{bbm}
\usepackage{dsfont}
\usepackage{xspace}
\usepackage{cancel}
\usepackage{mathtools}
\usepackage{amsthm}
\makeatletter
\@ifundefined{theorem}{\newtheorem{theorem}{Theorem}}{}
\makeatother
\title[Yet Another Distributional Bellman Equation]{Yet Another Distributional Bellman Equation (Theorem 5.2)}
\author{Nicole B\"auerle, Tamara G\"oll, Anna Ja\'skiewicz}
\date{Source: arXiv:2505.21098v2 (CC BY 4.0)}
\begin{document}
\maketitle

\noindent\emph{Source and licence.} Yet Another Distributional Bellman Equation. Version arXiv:2505.21098v2 (CC BY 4.0), distributed under the Creative Commons Attribution 4.0 International licence, as recorded in the arXiv licence metadata for this exact version and shown on the abstract page https://arxiv.org/abs/2505.21098v2. Full rights notice: licenses/arxiv-2505.21098-CC-BY-4.0.txt.\par\smallskip

\noindent\textit{Editorial scope.} Counted unit: the Theorem printed as Theorem 5.2 of the article (source0.tex lines 1027--1040, source label \texttt{theo:OT}), delivered together with the article's own complete original proof of it (source0.tex lines 1042--1080, one \texttt{proof} environment of 39 lines). No mathematics is written, repaired or re-derived: statement and proof are the article's own text line for line. Required context retained uncounted: none. Every definition, display and label of those lines is retained unchanged. Omitted from this package and not redistributed: 1--1026, 1041--1041, 1081--1208. Nothing inside a retained excerpt is omitted or altered: every retained body is delivered exactly as the article's lines are written, so no omission receipt of any kind accompanies this package. Citations: the retained excerpts carry no citation command at all, so no bibliography entry of the article is transcribed and no reference is redistributed. Reference adaptation: none was needed and none was made; every reference inside the delivered unit resolves inside it, so no unresolved reference or citation is produced. Printed slips kept literal: no printed word, symbol, hypothesis or number of the counted statement or of its proof was altered. Layout and class: the article's own class is \texttt{amsart} with packages including amsaddr, amssymb, hyperref, graphicx, listings, multirow, placeins, color, subfigure, lscape, dsfont, amsmath, amsfonts, tikz; the wrapper is the lane's pinned \texttt{amsart} editorial wrapper at 10pt on A4, which restates the theorem-like environments the retained text needs and adds the article's own macro definitions that the retained text needs (none), with extra wrapper packages (bm, bbm, dsfont, xspace, cancel, mathtools). The delivered numbering is the unit's own. Assets: no retained span contains a figure environment, a graphics-inclusion command, a PDF-page inclusion, a table or a TikZ picture; no image file is copied into this package, the package declares no asset dependency and no image file is redistributed with it. Licence: version arXiv:2505.21098v2 of this article is distributed under the Creative Commons Attribution 4.0 International licence, as recorded in the arXiv licence metadata for this exact version and shown on the abstract page https://arxiv.org/abs/2505.21098v2. A full rights notice is delivered at licenses/arxiv-2505.21098-CC-BY-4.0.txt. Characters: every retained excerpt is pure ASCII, so the delivered engine, XeLaTeX with its default Latin Modern text font, reports no missing character.
\begin{theorem}\label{theo:OT}
The following statements hold  for the optimal policy:
\begin{itemize}
\item[a)] If at time point $n$  not all mass 
at $x$ is moved, then $x$w ill only receive 
mass at later time points and stay a sink.
%, which means that  $F_m(x)\ge F_n(x)(1-\bar a_n^1(x)-\bar a_n^2(x))$for $m\ge n+1.$
%$x$ stay a sink. 
Formally, for any $x\in E$  
$$ \bar a_n^1(x)+\bar a_n^2(x)< 1 \quad \Rightarrow\quad \bar a_m^1(x)+\bar a_m^2(x)=0,\, m\ge  n+1. $$
\item[b)]  A mass which has already been moved 
will keep its direction of moving.
    \end{itemize}
\end{theorem}

\begin{proof} Part a): 
Assume that $\bar a_n^1(x)+\bar a_n^2(x)< 1$
but $\bar a_m^1(x)+\bar a_m^2(x)>0$
for some $m\ge n+1.$ Note that the cost from the coordinate $x$ equals
$$c_nF_n(x)(\bar a_n^1(x)+\bar a_n^2(x))$$
and $F_{n+1}(x)\ge F_n(x)(1-\bar a_n^1(x)-\bar a_n^2(x)).$   Let $m\ge n+1$ be the smallest  number for which
$a_m^1(x)+a^2_m(x)>0,$ which implies that $F_m(x)\ge F_{n+1}(x).$ 
Suppose that it is optimal in period $m$ is to transfer at least the mass $F_n(x)(\alpha_1\bar a_n^1(x)+\alpha_2\bar a_n^2(x)),$ where $\alpha_1,\alpha_2\in [0,1)$ are numbers such that 
$$0<\alpha_1\bar a_n^1(x)+\alpha_2\bar a_n^2(x)\le 1-
\bar a_n^1(x)-\bar a_n^2(x).$$
The cost from $x$ in period $m$ is 
$$c_mF_m(x)(\bar a_m^1(x)+\bar a_m^2(x))\ge
c_m F_n(x)(\alpha_1\bar a_n^1(x)+\alpha_2\bar a_n^2(x)).
$$
But if the mass were transported at $n,$ then the cost would be 
 $$c_n F_n(x)(\alpha_1 \bar a_n^1(x)+\alpha_2 \bar a_n^2(x))\le c_m  F_n(x)(\alpha_1\bar a_n^1(x)+\alpha_2\bar a_n^2(x)).$$
 Hence, it is not optimal  to transfer the mass 
 $F_n(x)(\alpha_1 \bar a_n^1(x)+\alpha_2 \bar a_n^2(x))$ at period $m.$ 
 It is cheaper to transfer it at period $n$. Then, 
 $\bar a^1_n(x)$ and $\bar a^2_n(x)$ would not be optimal.
 Therefore, it must hold $a_m^1(x)=a^2_m(x)=0$ for 
 $m\ge n+1.$

 Even if $x$ receives new mass at later time 
points this will not be moved, because without 
loss of generality we can assume that older 
mass is moved first.



%Part a): Let $(F_n(1),\ldots,F_n(K))$ be the 
%state at time point $n$ in the lifted MDP. Suppose first for some $x\in E$ we have  $ \bar a_n^1(x)=\bar a_n^2(x)=0,$ but  $\bar a_{n+1} ^1(x)+\bar a_{n+1}^2(x)>0.$  Then, set $\hat a_n^i(x):=\alpha \bar a_{n+1}^i(x), i=1,2$ for $\alpha = \min\{1,\alpha^*\}$ where $\alpha^*$ is such that $\alpha^* (\bar a_{n+1}^1(x)+\bar a_{n+1}^2(x)) =F_n(x)$ and $\hat a_{n+1}^i(x):= (1-\alpha) \bar a_{n+1}^i(x)$ and  define in the same manner $\hat a_n^i(y),\hat a_{n+1}^i(y)$ for $y\neq x$. Using This yields the same state $F_{n+2}$ and results in less transportation cost, since $c_n \le c_{n+1}$.Note that this argument also applies when $a_n^1(x)=a_n^2(x)>0$, but $a_n^1(x)=a_n^2(x) < F_n(x).$ The remaining mass $F_n(x) -(a_n^1(x)+a_n^2(x))$ will stay at $x$ forever. 


Part b): Moving a mass in one direction and 
back at a later time results in the same 
distribution and just produces costs. Hence, 
this cannot be optimal.
\end{proof}

\end{document}
